\documentclass[11pt,a4paper]{article}

\usepackage[utf8]{inputenc}

\usepackage[T1]{fontenc}

\usepackage[polish]{babel}

\usepackage{lmodern}

\usepackage{geometry}

\geometry{a4paper,margin=27mm}

\usepackage{microtype}

\usepackage[hidelinks]{hyperref}

% Portfolio: MSC-AR-P08; source: prospective ARC-XR research scope.

\title{Adaptacja poziomu wyzwania XR a transfer poprawnego wykonania zadań}

\author{Mateusz Pomianek}

\date{Wersja 0.3.0, 21 listopada 2026}

\begin{document}

\maketitle

\begin{abstract}

Przedmiotem opracowania jest Adaptacja poziomu wyzwania XR a transfer poprawnego wykonania zadań. Czy adaptacja trudności w XR poprawia transfer wykonania do nowego zadania? Projekt obejmuje stały oraz adaptowany poziom wyzwania podczas treningu, a następnie niezależny, odroczony test umiejętności bez adaptacji. Głównym punktem końcowym jest poprawność w zadaniu transferowym po odroczeniu, a nie tylko wynik treningu. Opis ma charakter wstępny i nie przedstawia niezweryfikowanych wyników.

\end{abstract}

\section{Wprowadzenie}

Rozwój adaptacyjnych systemów XR stwarza możliwość badania interakcji człowieka z informacją przestrzenną, ale samo zastosowanie AR, VR lub modeli sztucznej inteligencji nie przesądza o przewadze poznawczej czy ergonomicznej. W rozpatrywanym zagadnieniu, Adaptacja poziomu wyzwania XR a transfer poprawnego wykonania zadań, kluczowe jest rozdzielenie właściwości systemu od zmienności zachowania użytkownika. Porównanie należy zaprojektować tak, aby mierzyć jednoznaczne zjawisko, a nie tylko ogólną ocenę technologii.

\section{Cel pracy i pytanie badawcze}

Czy adaptacja trudności w XR poprawia transfer wykonania do nowego zadania? Projekt wymaga krytycznego przeglądu aktualnego stanu badań oraz zdefiniowania nowości w stosunku do innych tematów portfolio XR.

\section{Zakres i odrębny wkład naukowy}
Transfer sprawności do nowego, odroczonego zadania po adaptowanym versus statycznym poziomie wyzwania; nie ocena wyniku podczas treningu ani deklarowanego flow. Test trwałego transferu poznawczo-motorycznego, nie ogólna poprawa zaangażowania adaptacyjnym XR.
\section{Hipotezy badawcze}
\paragraph{H1.} Adaptacja trudności poprawia poprawność w niezależnym zadaniu transferowym względem stałego treningu przy wyrównanej ekspozycji.
\paragraph{H2.} Efekt transferu pozostaje dodatni przy kontrolowaniu liczby prób i czasu treningu, a nie tylko przy wyższej satysfakcji.
\paragraph{H3.} Różnica jest wykrywalna po odroczonym teście bez aktywnego adaptatora, z niezależną wersją zadania.
\section{Projekt badania i plan eksperymentów}
\subsection{A. Operacjonalizacja i kwalifikacja danych}
Transfer sprawności do nowego, odroczonego zadania po adaptowanym versus statycznym poziomie wyzwania; nie ocena wyniku podczas treningu ani deklarowanego flow. Jednostka wnioskowania: uczestnik × nowy scenariusz transferowy. W protokole zamrozić definicję i mianownik wyniku oraz źródło prawdy odniesienia, zgodność znaczników czasu, prawa do danych i zasady rekrutacji.
\subsection{B. Eksperyment porównawczy}
Zrównoważyć trudność bazową, losować przydział i identycznie ograniczyć czas treningu; następnie zmierzyć wykonanie niećwiczonej sekwencji. Estymand pierwszorzędowy: różnica prawdopodobieństwa poprawnego wykonania niećwiczonego zadania po identycznym czasie treningu. Przed testem zdefiniować randomizację, próbę, punkt końcowy i warunki wykluczenia bez selekcji wyniku post hoc.
\subsection{C. Kontrole i analiza odporności}
Oddzielić efekt uczenia konkretnych elementów od przeniesienia strategii, zastosować nową geometrię i reguły zadania. Raportować czas, błędy i niedokończone transfery bez zastępowania rezultatu głównego wynikiem ćwiczenia. W rejestrze uwzględnić braki, wstrzymania, opóźnienia i niepełne próby zamiast usuwać je z analizy.
\subsection{D. Walidacja na niezależnym materiale}
Weryfikować na niezależnych sesjach lub osobach, bez dostępu modelu do danych końcowych. Model dwumianowy z uczestnikiem i zadaniem jako efektami, analiza intencji leczenia i klastrowy przedział różnicy ryzyka. Dokumentować wyniki niesprzyjające hipotezie oraz rzeczywiste granice przeniesienia.
\section{Estymand, analiza statystyczna i decyzje}
różnica prawdopodobieństwa poprawnego wykonania niećwiczonego zadania po identycznym czasie treningu. Jednostka: uczestnik × nowy scenariusz transferowy. Model dwumianowy z uczestnikiem i zadaniem jako efektami, analiza intencji leczenia i klastrowy przedział różnicy ryzyka. H1 jest kandydatem hipotezy potwierdzającej, a H2/H3 pozostają eksploracyjne do formalnego planu wielokrotnych porównań. Liczebność uczestników, margines, próg i moc lub precyzja wymagają pilotażu oraz zamrożenia przed testem końcowym.
\section{Stan badań i granica nowości}
\begin{itemize}
\item XR personalization methods exist; a claim needs a defined multi-session or transfer contrast. \cite{xrprior1}.
\item Learning-load effects do not imply delayed knowledge retention or skill transfer. \cite{xrprior2}.
\end{itemize}
Test trwałego transferu poznawczo-motorycznego, nie ogólna poprawa zaangażowania adaptacyjnym XR. Odrębność od wcześniejszych publikacji i innych tematów programu należy sprawdzić przed rejestracją badania.
\section{Warunki bezpieczeństwa, sprzętu i zgodności}
Jest to plan prospektywny, nie dowód wykonania eksperymentu lub zatwierdzonej etyki. Udział osób, rejestrowanie biometrii i interwencje wymagają uprzedniej procedury etycznej i świadomej zgody. Nie wnioskować o emocjach ani zdrowiu z samych sygnałów. Standardowy Quest 3 nie ma okulometrii HMD, a Tobii Eye Tracker 5 służy do stanowiska ekranowego; każde połączenie zweryfikować przed protokołem.
\section{Materiał i metody}

Projekt obejmuje stały oraz adaptowany poziom wyzwania podczas treningu, a następnie niezależny, odroczony test umiejętności bez adaptacji. Wymagane jest zdefiniowanie jednostki analizy, zasad doboru przypadków, wersji scenariusza i źródeł pomiaru. Rozkład warunków oraz reguły wykluczenia prób należy określić przed pomiarem. W badaniach z udziałem ludzi niezbędna jest uprzednia właściwa ocena etyczna i świadoma zgoda. W testach wyłącznie symulacyjnych należy oddzielić prawdę symulacyjną od obserwacji empirycznej.

\subsection{Wynik główny i kryteria jakości}

Poprawność w zadaniu transferowym po odroczeniu, a nie tylko wynik treningu. Pomiar powinien zachowywać jawne jednostki, znaczniki czasu i reguły oceny poprawności. Przy modelach predykcyjnych stosować rozłączne partycje oceny; przy eksperymentach interwencyjnych oddzielić skuteczność predykcji od przyczynowego wpływu interfejsu.

\section{Plan analizy}

Porównanie zmiany od pomiaru bazowego do testu odroczonego z kontrolą różnic początkowych i efektów kolejności. W wynikach należy raportować wielkość efektu lub błąd estymacji z przedziałem niepewności, liczbą obserwacji i odsetkiem odrzuconych prób. Analizy dodatkowe wymagają jasnego oznaczenia jako eksploracyjne. Nie wybierać progów decyzyjnych ani modelu na podstawie obserwacji ze zbioru testowego.

\section{Wyniki}

W obecnej wersji nie raportuje się wyników pomiarów, tabel statystycznych ani parametrów wyznaczonych z nieudokumentowanych danych. To kierunek badawczy, a nie zatwierdzony protokół badania ani zgłoszony artykuł.

\section{Dyskusja}

Interpretacja przyszłych wyników musi uwzględnić alternatywne mechanizmy i zależność wniosków od jakości danych, warunków zadania oraz ograniczeń sprzętu. Wartość potencjalnego rozwiązania zależy nie tylko od trafności modelu, lecz także od jego opóźnienia, stabilności i wpływu na zachowanie człowieka. Wyniki testów syntetycznych nie mogą być automatycznie uogólniane na populację uczestników.

\section{Ograniczenia badania}

Trening w trudniejszych warunkach może zmniejszać bieżący komfort, ale niekoniecznie pogarszać uczenie. Wynik wykonania zadania, nie subiektywne doświadczenie flow. Żadna część szkicu nie dowodzi przeprowadzenia zbierania danych, uzyskania zgód etycznych ani gotowości manuskryptu do zgłoszenia.

\section{Wnioski}

Przedstawiono problem, kontrast i sposób jego planowanej oceny. Rozstrzygnięcie postawionego pytania wymaga przygotowania protokołu, zweryfikowanych źródeł oraz poprawnie udokumentowanych wyników. Bez tych materiałów nie można sformułować konkluzji empirycznych.

\begin{thebibliography}{9}
\bibitem{xrprior1} Adapting to the User: A Systematic Review of Personalized Interaction in VR (2025), arXiv:2510.13123. \url{https://arxiv.org/abs/2510.13123}.
\bibitem{xrprior2} Can augmented reality technology reduce learners' cognitive load? A meta-analysis (2026), Smart Learning Environments. \url{https://doi.org/10.1186/s40561-025-00429-7}.
\end{thebibliography}
\end{document}

